% The reverse-complement quotient is a signed graph
% Build: pdflatex rc_quotient_signed (x3)
\documentclass[11pt]{article}
\usepackage{lmodern}
\usepackage[T1]{fontenc}
\usepackage{amsmath,amssymb,amsthm,mathtools}
\usepackage[margin=1.1in]{geometry}
\usepackage{booktabs}
\usepackage{microtype}
\usepackage[hidelinks]{hyperref}

\theoremstyle{plain}
\newtheorem{theorem}{Theorem}
\newtheorem{proposition}[theorem]{Proposition}
\newtheorem{lemma}[theorem]{Lemma}
\newtheorem{corollary}[theorem]{Corollary}
\theoremstyle{definition}
\newtheorem{definition}[theorem]{Definition}
\theoremstyle{remark}
\newtheorem{remark}[theorem]{Remark}

\newcommand{\Z}{\mathbb{Z}}
\newcommand{\Q}{\mathbb{Q}}
\newcommand{\K}{\mathcal{K}}
\newcommand{\coker}{\operatorname{coker}}
\newcommand{\im}{\operatorname{im}}
\newcommand{\mult}{\operatorname{mult}}
\newcommand{\rank}{\operatorname{rank}}
\newcommand{\Sig}{\Sigma}

\title{\bfseries The reverse-complement quotient is a signed graph:\\
what Zaslavsky's theory gives, and what it cannot}
\author{Ruqing Chen\\[2pt]
\normalsize GUT Geoservice Inc., Rivi\`ere-Beaudette, Qu\'ebec J0P 1R0, Canada\\
\normalsize\texttt{ruqing@hotmail.com}}
\date{20 September 2026}

\begin{document}
\maketitle

\begin{abstract}
\noindent
The companion note built the reverse-complement double cover $D_k(S)$ of a de Bruijn
multidigraph and found that the involution $\rho(v)=\bar v$ carries $D_k$ to its
\emph{opposite} digraph, so that $\rho$ induces no endomorphism of $\K(D_k)$ and the cover
does not descend. The obvious repair is to pass to the quotient $\Sig_k=D_k/\rho$ and to use
Zaslavsky's signed Laplacian there. Half of that works, and the half that fails is the
interesting half. \emph{Works:} $\rho$ is free exactly when no repeated $k$-mer is its own
reverse complement, which is automatic for odd $k$ and genuinely fails for even $k$
($\phi$X174 has three palindromic repeated $10$-mers and one palindromic repeated $12$-mer);
every bundle of $\Sig_k$ is monochromatic, so $\Sig_k$ carries honest $r$-bundles; $\Sig_k$
is \emph{unbalanced} precisely when $S$ carries an inverted repeat of length $\ge k$, which
is Zaslavsky balance read through the connectivity criterion of the companion note; and the
Reiner--Tseng sequence
$0\to\K(\Sig_k)\to\K(|D_k|)\to\K(|\Sig_k|)\to 0$
holds on the nose, with a non-split $2$-primary part.
\emph{Fails, twice.} First, if $P$ is the permutation matrix of $\rho$ then
$P\Delta(D_k)P=\Delta(D_k)^{\mathsf T}$ and never $\Delta(D_k)$: the deck involution is an
\emph{anti}-automorphism of the Laplacian, the $\pm1$ eigenlattices of $P$ are not
$\Delta$-stable, and the \emph{directed} sandpile group --- the one the BEST theorem counts
reconstructions with --- lies in no covering sequence. What it gives instead is a perfect
\emph{symmetric} linking form $\K(D_k)\times\K(D_k)\to\Q/\Z$, so the relevant notion is a
metabolizer rather than a Lagrangian; and since $|\K(D_{11})|=25\,318\,912$ is not a square,
$\phi$X174 has none. Second, a signed $r$-bundle chain does
not contract: the presentation column at the chain's endpoint differs between $\Sig$ and
$\Sig/C$ by exactly $s\,r\,d_0$, the series law, so the splitting theorem that makes the
genome-scale computation feasible on one strand has no bidirected analogue. What does survive
is the $r$-rank. The seven ribosomal operons of \emph{E.\ coli} K-12 merge into genuine
$7$-bundles of $\Sig_k$ whose longest block spans $728$\,bp at $k=31$, $51$ and $75$ alike ---
the same block the companion note matched one-to-one onto the seven operons --- and the only
trace the quotient keeps of them is the $\mathbb{F}_7$-nullity of the signed Laplacian.
\end{abstract}

\section{What was asked, and what is true}

The single-strand theory of these notes assigns to a sequence $S$ and a word length $k$ the
de Bruijn multidigraph $G_k(S)$, whose vertices are the repeated $k$-mers and whose arcs are
the $(k{+}1)$-mer occurrences, and reads the reconstruction count off the critical group
$\K(G_k)$ through the BEST theorem. Double-strandedness was handled in the companion note by
passing to the \emph{double cover} $D_k(S)$: the multidigraph whose arcs are the
$(k{+}1)$-mer occurrences in $S$ \emph{and} in $\bar S$. That graph is balanced, the unitig
lemma and the splitting theorem apply to it verbatim, and the seven ribosomal operons of
\emph{E.\ coli} --- five on one strand, two on the other --- finally appear as one
seven-copy family instead of a five-copy and a two-copy family.

What the companion note could not do was quotient. The map $\rho(v)=\bar v$ satisfies
\[
  \rho\colon D_k \longrightarrow D_k^{\mathrm{op}},\qquad
  \rho(u\to v)=\bar v\to\bar u ,
\]
so it reverses arcs; the quotient of a digraph by a direction-reversing involution is not a
digraph, and $\rho$ induces a map $\K(D_k)\to\K(D_k^{\mathrm{op}})\cong\K(D_k)^\vee$ rather
than an endomorphism. The standard repair in combinatorics is to forget the directions, to
record instead a \emph{sign} on each edge of the quotient, and to work with Zaslavsky's
signed Laplacian $\Delta_\Sig=D-A_\sigma$ \cite{Zas82}. This note carries that out, and
reports both what it yields and the two places where it provably cannot yield what was
wanted.

The one-paragraph summary is this. The quotient is a bona fide signed graph, its critical
group is a switching invariant of the genome, its balance has an exact biological meaning,
and the integral covering theory of Reiner and Tseng \cite{RT14} applies to it word for word.
But that theory computes an invariant of the \emph{undirected shadow} $|D_k|$, which by the
matrix-tree theorem counts spanning trees; the quantity of interest counts Eulerian circuits
and lives on the digraph, and Theorem~\ref{thm:anti} says the digraph admits no such theory.
Separately, Theorem~\ref{thm:series} says the chain contraction that reduces a bacterial
genome from tens of thousands of vertices to a few hundred has no signed analogue. Together
these explain why the ``unified bidirected sandpile group'' that one would like to write down
does not exist, and what the correct residue of the seven operons is instead.

Throughout, $S$ is a circular sequence of length $N$ over $\{A,C,G,T\}$, $\bar S$ its reverse
complement, and $\bar w$ the reverse complement of a word $w$. All computations are exact
integer arithmetic; the script \texttt{bidirected\_sandpile.py} reproduces every number below
and prints a seven-line pass/fail battery.

\section{The quotient, and when it exists}

\begin{definition}
$D_k=D_k(S)$ is the multidigraph of the companion note: one vertex per $k$-mer occurring at
least twice in the concatenation $S\bar S$, one arc $u\to v$ per occurrence of a
$(k{+}1)$-mer $w$ with prefix $u$ and suffix $v$ in $S$ or in $\bar S$. The involution
$\rho$ sends the vertex $w$ to $\bar w$ and the arc $u\to v$ to $\bar v\to\bar u$.
\end{definition}

The first question is whether $\rho$ acts freely on vertices, because the whole signed-graph
dictionary presupposes a free $\Z/2$ action.

\begin{proposition}\label{prop:parity}
$\rho$ fixes a vertex of $D_k$ if and only if some repeated $k$-mer $w$ satisfies
$\bar w=w$. For odd $k$ no such word exists over any alphabet with a fixed-point-free
complement, so $\rho$ is free. For even $k$ such words exist and do occur.
\end{proposition}

\begin{proof}
A word $w$ of odd length $k$ with $\bar w=w$ would have its middle letter equal to its own
complement, and complementation on $\{A,C,G,T\}$ has no fixed point. For even $k$ the
condition $\bar w=w$ is the condition that $w$ be a palindromic restriction-site-like word,
and these are common.
\end{proof}

\begin{table}[t]
\centering
\begin{tabular}{rlrrl}
\toprule
$k$ & parity & $|V(D_k)|$ & $\rho$-fixed & $A_\sigma$ symmetric\\
\midrule
 9 & odd  & 396 & 0 & yes\\
10 & even & 125 & 3 & \textbf{no}\\
11 & odd  &  30 & 0 & yes\\
12 & even &   7 & 1 & \textbf{no}\\
\bottomrule
\end{tabular}
\caption{$\phi$X174 (NC\_001422.1). Palindromic repeated $k$-mers exist at $k=10$ and
$k=12$, and at exactly those $k$ the quotient adjacency matrix fails to be symmetric, so no
signed graph is defined. Check (P).}
\label{tab:parity}
\end{table}

Table~\ref{tab:parity} makes the point concretely and shows that the obstruction is not
hypothetical: it is precisely at the even $k$ with palindromic repeats that the construction
below breaks, and it breaks visibly, in the symmetry of the adjacency matrix. From here on
$k$ is odd, or even with no palindromic repeated $k$-mer.

\begin{remark}[What the fixed points are, topologically]\label{rem:halfedge}
The failure at even $k$ is not a technicality but a change of category. A vertex $w=\bar w$
is a branch point of the projection $D_k\to D_k/\rho$, so the quotient is not a covering
space and the fibre over $w$ is a single point rather than two. In Zaslavsky's language the
image acquires a \emph{half-edge} --- an edge with one endpoint --- at $w$: an arc $u\to w$
and its partner $w\to\bar u$ descend to the same edge of the quotient but attach to $[w]$
with opposite end-signs, since $\varepsilon(w)=+1$ forces the two signs
$\varepsilon(u)\varepsilon(w)$ and $\varepsilon(w)\varepsilon(\bar u)$ to differ. That is
exactly why the quotient adjacency matrix in Table~\ref{tab:parity} fails to be symmetric:
$A_\sigma$ records one sign per edge, and a half-edge has two. Reiner and Tseng do treat
signed graphs with half-edges \cite[\S9]{RT14}, so the object is not outside the theory, but
its critical group is no longer the one computed below, and the double cover it belongs to is
branched. Biologically the fixed points are the palindromic restriction-site-like words, so
the branch locus of the quotient is precisely the repeated restriction sites of the genome.
\end{remark}

\begin{definition}\label{def:sigma}
Choose a representative of each $\rho$-orbit $\{v,\bar v\}$ and set $\varepsilon(v)=+1$ if
$v$ is the chosen one and $-1$ otherwise. The signed graph $\Sig_k=D_k/\rho$ has one vertex
per $\rho$-orbit of vertices, one edge per $\rho$-orbit of arcs, and the edge coming from
$u\to v$ carries the sign $\varepsilon(u)\varepsilon(v)$.
\end{definition}

The arcs do fall into orbits of size two, and here no parity hypothesis is needed. An arc of
$D_k$ is an \emph{occurrence}: a position in $S$ or a position in $\bar S$. The involution
$\rho$ carries occurrences in $S$ to occurrences in $\bar S$ and back, and no occurrence lies
in both, so $\rho$ is free on arcs outright. The script nevertheless constructs the pairing
explicitly and fails loudly if any arc lacks a partner.

Two facts make Definition~\ref{def:sigma} usable.

\begin{lemma}\label{lem:mono}
Every bundle of parallel edges of $\Sig_k$ is monochromatic.
\end{lemma}

\begin{proof}
The sign $\varepsilon(u)\varepsilon(v)$ depends only on the endpoints of the arc, not on the
arc. All $r$ arcs $u\to v$ therefore descend to $r$ edges of the same sign.
\end{proof}

This is what licenses the phrase ``$r$-bundle'' below, and it is not automatic for a general
signed multigraph.

\begin{lemma}\label{lem:switch}
Changing the choice of representatives is a Zaslavsky switching, and
$\Delta_{\Sig_k}$ changes by conjugation by a diagonal $\pm1$ matrix. Hence $\coker
\Delta_{\Sig_k}$, and with it the critical group of Section~\ref{sec:rt}, depends only on
$(S,k)$.
\end{lemma}

Replacing the representative at one orbit multiplies $\varepsilon$ by $-1$ there, which
multiplies the corresponding row and column of $A_\sigma$ by $-1$ and leaves $D$ alone. The
script verifies invariance directly on eight random representative sets at $k=9$ and $k=11$,
obtaining for instance
\[
  \K(\Sig_{11})=(\Z/2)^2\oplus\Z/4204481
\]
every time. Check (W).

\section{Balance is an inverted repeat}

\begin{theorem}\label{thm:balance}
$\Sig_k$ is unbalanced if and only if $S$ carries an inverted repeat of length at least $k$.
\end{theorem}

\begin{proof}
A signed graph is balanced if and only if its signed double cover is disconnected; for the
$\Z/2$-voltage double cover this is Reiner and Tseng, Propositions 11.4 and 11.5
\cite{RT14}. The underlying multigraph of $D_k$ \emph{is} that double cover, by
Definition~\ref{def:sigma}. The companion note proves that $D_k$ is connected exactly when
$S$ carries an inverted repeat of length at least $k$, the two components otherwise being
the plus-strand and minus-strand halves.
\end{proof}

For $\phi$X174 the longest inverted repeat has length $12$, realised by
\texttt{ATTAAGCTCATT} at position $441$ against its reverse complement at $1448$. The script
confirms the trichotomy at $k=9,11,13,15$: unbalanced and connected at $9$ and $11$,
balanced and disconnected at $13$ and $15$. Check (B).

Theorem~\ref{thm:balance} is worth stating because it converts a purely combinatorial
dichotomy into a statement about the sequence. A genome with no inverted repeat as long as
$k$ gives a balanced $\Sig_k$; balanced means switching-equivalent to all-positive, which
means the signs carry no information, which means the whole construction has collapsed back
to the single-strand one. The sign structure is \emph{exactly} the record of inverted
repeats, nothing more and nothing less.

\section{The covering theory that does exist}\label{sec:rt}

Reiner and Tseng define the critical group of a signed graph not as $\coker\Delta_\Sig$ but
as $\K(\Sig)=\im\partial/\im\Delta_\Sig$, where $\partial$ is Zaslavsky's signed incidence
matrix \cite[eq.~(1.3), \S9]{RT14}. For a connected unbalanced $\Sig$ one has
$\im\partial=\{x\in\Z^V:\sum_v x_v\ \text{even}\}$, an index-two sublattice of $\Z^V$, so
$\K(\Sig)$ is an index-two subgroup of $\coker\Delta_\Sig$; ignoring the distinction doubles
every order and corrupts the $2$-primary part. (The competing convention of Cho, Dochtermann,
Inagaki, Oh, Snustad and Zacovic \cite{CDIOSZ} deletes a sink row and column instead, which
is the chip-firing-natural choice but a different group.) We use the Reiner--Tseng
convention, computed from the explicit basis $2e_0,\ e_0+e_1,\ \dots,\ e_0+e_{n-1}$ of
$\im\partial$.

\begin{theorem}[Reiner--Tseng \cite{RT14}, Thm.~1.2, verified here]\label{thm:rt}
For every odd $k$ with $\Sig_k$ unbalanced there is a short exact sequence
\[
  0\longrightarrow \K(\Sig_k)\longrightarrow \K(|D_k|)\longrightarrow \K(|\Sig_k|)
  \longrightarrow 0 ,
\]
split on $p$-primary components for odd $p$ and in general \emph{not} split at $p=2$.
\end{theorem}

The $\phi$X174 instance at $k=11$ is a textbook illustration:
\[
\begin{aligned}
  \K(\Sig_{11})   &= (\Z/2)^2\oplus\Z/4204481, &&\text{order } 16\,817\,924,\\
  \K(|D_{11}|)    &= (\Z/2)^2\oplus\Z/4\oplus\Z/23\oplus\Z/12757\oplus\Z/4204481,
                  &&\text{order } 19\,738\,255\,595\,056,\\
  \K(|\Sig_{11}|) &= (\Z/2)^2\oplus\Z/23\oplus\Z/12757, &&\text{order } 1\,173\,644 .
\end{aligned}
\]
The orders multiply exactly. The odd parts split visibly: $4204481$ comes from the signed
graph, $23\cdot 12757$ from the base. The $2$-primary parts do \emph{not}: an extension of
$(\Z/2)^2$ by $(\Z/2)^2$ has produced $(\Z/2)^2\oplus\Z/4$, of the right order $16$ and the
wrong isomorphism type. This is precisely the failure Reiner and Tseng exhibit with
$0\to\Z_2\to\Z_4\to\Z_2\to 0$, arising here from a real bacteriophage. Check (R).

\section{The deck involution is an anti-automorphism}

Theorem~\ref{thm:rt} is a statement about $|D_k|$, the undirected shadow. The quantity the
single-strand theory actually computes is $\K(D_k)$, the sandpile group of the
\emph{digraph}, because that is what the BEST theorem multiplies by
$\prod_v(d^+(v)-1)!$ to count reconstructions. The following says the two are unrelated, and
says why.

\begin{theorem}\label{thm:anti}
Let $\Delta=\Delta(D_k)$ be the out-degree Laplacian of $D_k$ and $P$ the permutation matrix
of $\rho$. Then
\[
  P\,\Delta\,P=\Delta^{\mathsf T},
\]
and $P\Delta P\ne\Delta$ whenever $D_k$ has an arc that is not part of a $\rho$-symmetric
digon. Consequently $\Delta$ does not commute with $P$, the $\pm1$ eigenlattices of $P$ are
not $\Delta$-stable, and there is no decomposition of $\Delta$ over $\Z$, or over
$\Z[1/2]$, or over any ring, into a part on the invariants and a part on the anti-invariants.
\end{theorem}

\begin{proof}
$\rho$ is a bijection from the arc set of $D_k$ to itself carrying $u\to v$ to
$\bar v\to\bar u$. Hence the adjacency matrix satisfies $A_{\bar v\bar u}=A_{uv}$, that is
$PAP=A^{\mathsf T}$, and $d^+(\bar v)=d^-(v)=d^+(v)$ since $D_k$ is balanced, so $PDP=D$.
Subtracting gives $P\Delta P=\Delta^{\mathsf T}$. If $P\Delta P$ were also $\Delta$ then
$\Delta=\Delta^{\mathsf T}$, forcing $A=A^{\mathsf T}$, i.e.\ every arc paired with a reverse
arc of the same multiplicity.

For the consequence: a decomposition along the eigenlattices of $P$ requires
$\Delta P=P\Delta$. Here $\Delta P=P\Delta^{\mathsf T}$, and $\Delta^{\mathsf T}\ne\Delta$.
\end{proof}

This is the structural reason the companion note's obstruction was not an accident of its
proof. The classical covering-graph factorisation --- of Bilu and Linial \cite{BL06}, in the
form used for signed graphs --- decomposes the Laplacian of a double cover into an unsigned
and a signed part on the $\pm1$ eigenspaces of the deck transformation. That decomposition
needs the deck transformation to \emph{commute} with the Laplacian. Here it conjugates the
Laplacian to its transpose instead, and the two eigenlattices are exchanged rather than
preserved.

The numbers are unambiguous. For $\phi$X174 at $k=11$,
\[
  |\K(D_{11})| = 25\,318\,912 ,\qquad
  |\K(|D_{11}|)| = 19\,738\,255\,595\,056 ,
\]
and the second is not a multiple of the first; at $k=9$ the two orders are $124$- and
$205$-digit integers, again with no divisibility. The directed group is not an extension of
anything by $\K(\Sig_k)$ and is not a quotient of anything by it. Check (N).

\begin{remark}[What $\rho$ gives instead: a symmetric linking form]\label{rem:pairing}
The honest formulation is that $\rho$ makes $D_k$ into a digraph with a \emph{duality}, not a
digraph with a symmetry. Bowen--Franks duality (Paper VIII, Lem.~1.5 of the main series) says
$\coker M$ and $\coker M^{\mathsf T}$ are canonically dual for a square integer matrix $M$
with finite cokernel, and Theorem~\ref{thm:anti} identifies $\Delta^{\mathsf T}$ with
$P\Delta P$. Composing gives a perfect pairing on the sandpile group itself:
\[
  \langle u,v\rangle=(Pu)^{\mathsf T}\Delta^{-\mathsf T}v \bmod \Z,
  \qquad \K(D_k)\times\K(D_k)\longrightarrow\Q/\Z .
\]
It is well defined --- if $u=\Delta^{\mathsf T}z$ then $Pu=P\Delta^{\mathsf T}P\,(Pz)=
\Delta(Pz)$, so $\langle u,v\rangle=(Pz)^{\mathsf T}v\in\Z$ --- and it is
\emph{symmetric}, not alternating: inverting $P\Delta P=\Delta^{\mathsf T}$ gives
$P\Delta^{-\mathsf T}=\Delta^{-1}P$, whence
$\langle u,v\rangle=u^{\mathsf T}P\Delta^{-\mathsf T}v=u^{\mathsf T}\Delta^{-1}Pv=
\langle v,u\rangle$. Verified exactly for $\phi$X174 at $k=11$ and $k=12$; check (D).

This matters for how one should look for a double-stranded count. A symmetric linking form
has no Lagrangians in the symplectic sense; the relevant notion is a \emph{metabolizer}, a
subgroup $L$ with $L=L^{\perp}$, and a metabolizer forces $|\K(D_k)|=|L|^{2}$. For $\phi$X174
that already fails: $|\K(D_{12})|=33$ and $|\K(D_{11})|=25\,318\,912$, neither a square. So
whatever counts double-stranded reconstructions, it is not the order of a self-orthogonal
subgroup of $\K(D_k)$, and the symmetry of the form --- not merely its existence --- is the
first thing a candidate theory has to respect.

One further negative is worth recording, because it removes the most natural-looking
formulation. The sandpile group is a torsor on the arborescences \emph{converging} on a fixed
root, and those are what the BEST theorem counts with. One would like to say that a
double-stranded count is a sum over the $\rho$-orbits of arborescences. There are none:
$\rho$ reverses arcs, so it carries an arborescence converging on $r$ to one
\emph{diverging} from $\rho(r)$, and the two families are disjoint. What $\rho$ gives at
tree level is therefore a \emph{bijection}
\[
  \rho\colon \{\text{arborescences converging on } r\}
  \;\xrightarrow{\ \sim\ }\;
  \{\text{arborescences diverging from } \rho(r)\},
\]
both families having $|\K(D_k)|$ elements --- the tree-level shadow of the duality, not an
action. Verified for $\phi$X174 at $k=12$ at three roots, where each family has $33$ members
and the image lands in the diverging family every time and in the converging one never;
check (D). A double-stranded theory must be built from that bijection.
\end{remark}

\section{Signed bundle chains do not contract}\label{sec:series}

The single-strand computation for a bacterial genome is only possible because of the
splitting theorem: a \emph{chain} $C=(c_1,\dots,c_s)$ of vertices each of which sends all
$r$ of its out-arcs to the next splits off a factor $(\Z/r)^{s-1}$ and can be contracted,
reducing \emph{E.\ coli} at $k=75$ from $28\,579$ vertices to a skeleton of $261$. The
question is whether this survives in $\Sig_k$. It does not, and the reason is elementary
enough to state exactly.

\begin{theorem}\label{thm:series}
Let $\Sig$ be a signed graph containing a chain $C=(c_1,\dots,c_s)$: vertices of degree
exactly $2r$, joined by $r$-bundles to $c_{t-1}$ and $c_{t+1}$, where $c_0=a$ and
$c_{s+1}=b$ lie outside $C$. Switch so that every chain bundle is positive, and write
$d_t=x_{c_t}-x_{c_{t+1}}$ for $t=0,\dots,s$ in $\coker\Delta_\Sig$, where $x_v$ is the class
of the basis vector at $v$. Then
\begin{enumerate}
\item[\rm(a)] the presentation columns at $c_1,\dots,c_s$ say exactly
      $r(d_{t-1}-d_t)=0$ for $t=1,\dots,s$; in particular the $s$ second differences
      $\delta_t=x_{c_{t-1}}-2x_{c_t}+x_{c_{t+1}}$ are killed by $r$, and
      $r d_0=r d_1=\dots=r d_s$;
\item[\rm(b)] the column at $a$ in $\Sig$ and the column at $a$ in the contraction $\Sig/C$
      --- which replaces $C$ by a single $r$-bundle $a\!-\!b$ of the product sign --- differ
      by exactly $s\,r\,d_0$.
\end{enumerate}
Consequently $\K(\Sig)\not\cong(\Z/r)^s\oplus\K(\Sig/C)$ in general.
\end{theorem}

\begin{proof}
(a) $\Delta_\Sig$ is symmetric and its column at $c_t$ has entry $2r$ at $c_t$ and $-r$ at
each of $c_{t-1},c_{t+1}$, all other entries zero. The relation it imposes is
$2r\,x_{c_t}-r\,x_{c_{t-1}}-r\,x_{c_{t+1}}=0$, which is $-r(d_{t-1}-d_t)=0$. Chaining the
$s$ relations gives $r d_0=\dots=r d_s$.

(b) The column at $a$ in $\Sig$ contains $-r\,x_{c_1}$ where the column at $a$ in $\Sig/C$
contains $-r\,x_b$; all other entries agree, since $\deg a$ is $r$ in both. The difference is
$r(x_b-x_{c_1})=-r\sum_{t=1}^s d_t$, and by (a) each $r d_t$ equals $r d_0$, so this is
$-s\,r\,d_0$.
\end{proof}

Contrast the directed case. There the column at $c_t$ is $r(x_{c_t}-x_{c_{t+1}})=r d_t$,
which kills each \emph{first} difference outright; the chain contributes $(\Z/r)^{s-1}$ as a
direct summand and contracts cleanly. In the signed case only second differences die. The
residual first difference $d_0$ is not torsion; it is coupled to the rest of the graph, and
the coupling is the series law of electrical networks --- $s{+}1$ bundles of conductance $r$
in series behave like one bundle of conductance $r/(s{+}1)$, which is not an integer
multiplicity, so there is no contracted signed graph to compare to.

\begin{table}[t]
\centering
\begin{tabular}{rrllrr}
\toprule
$r$ & $s$ & $\K(\Sig)$ & $\K(\Sig/C)$ & ratio & $r^s$\\
\midrule
 3 & 1 & $(\Z/3)^2\oplus\Z/8$              & $\Z/2\oplus\Z/9$          &   4.00 &   3\\
 3 & 2 & $\Z/2\oplus(\Z/3)^3\oplus\Z/8$    & $(\Z/2)^2\oplus\Z/9$      &  12.00 &   9\\
 3 & 3 & $(\Z/3)^3\oplus\Z/4\oplus\Z/9$    & $\Z/2\oplus\Z/9$          &  54.00 &  27\\
 5 & 2 & $(\Z/2)^2\oplus(\Z/5)^2\oplus\Z/17$ & $\Z/2\oplus\Z/4\oplus\Z/7$ & 30.36 &  25\\
 5 & 3 & $\Z/4\oplus(\Z/5)^3\oplus\Z/11$   & $\Z/2\oplus\Z/13$         & 211.54 & 125\\
 7 & 2 & $\Z/2\oplus(\Z/7)^2\oplus\Z/23$   & $\Z/2\oplus\Z/17$         &  66.29 &  49\\
 7 & 3 & $\Z/4\oplus(\Z/7)^3\oplus\Z/13$   & $\Z/2\oplus\Z/17$         & 524.59 & 343\\
11 & 2 & $\Z/2\oplus(\Z/11)^2\oplus\Z/64$  & $(\Z/2)^2\oplus\Z/29$     & 133.52 & 121\\
\bottomrule
\end{tabular}
\caption{An $r$-bundle chain of $s$ internal vertices grafted into a fixed unbalanced base.
The $r$-rank rises by exactly $s$ in all eight rows --- the $(\Z/r)^s$ of
Theorem~\ref{thm:series}(a) is visible --- but the order ratio never equals $r^s$, and is
not even an integer in five rows. Check (C).}
\label{tab:series}
\end{table}

Table~\ref{tab:series} separates the two halves of the theorem. What survives contraction is
the $r$-\emph{rank}: the number of invariant factors divisible by $r$ grows by exactly $s$,
in every case tested. What does not survive is the order and the isomorphism type. This is a
partial and negative data point for Question 5.5 of \cite{CDIOSZ}, which asks how
$\K(\Sig_\varphi)$ depends on $\K(|\Sig|)$ as the signature varies and observes only that the
number of invariant factors seems to be preserved.

\begin{corollary}\label{cor:nocompute}
The reduction that makes $\K(G_k)$ computable for a bacterial genome --- contract every
maximal chain, compute on the skeleton, multiply back --- has no analogue for $\K(\Sig_k)$.
\end{corollary}

\section{\emph{E.\ coli} K-12: what is left of the seven operons}

\begin{table}[t]
\centering
\begin{tabular}{rrrrrrrr}
\toprule
$k$ & $|V(D_k)|$ & $|V(\Sig_k)|$ & $\rho$-fixed & $7$-bundles & longest $7$-chain
& span (bp) & $\dim_{\mathbb{F}_7}\ker\Delta_{\Sig_k}$\\
\midrule
31 & 60\,548 & 30\,274 & 0 & 4919 & 698 & \textbf{728} & 7441\\
51 & 53\,380 & 26\,690 & 0 & 4352 & 678 & \textbf{728} & 6581\\
75 & 48\,498 & 24\,249 & 0 & 3814 & 654 & \textbf{728} & 5864\\
\bottomrule
\end{tabular}
\caption{\emph{E.\ coli} K-12 MG1655, NC\_000913.3, $4\,641\,652$\,bp, at the three odd word
lengths of the single-strand study. The longest $7$-bundle chain spans the same $728$
base pairs at all three $k$ --- $698+31-1=678+51-1=654+75-1$ --- which is the seven-copy
block the companion note matched one-to-one onto the seven ribosomal operons. Check (E).}
\label{tab:ecoli}
\end{table}

The construction does what was asked of it structurally. $\rho$ is free at all three odd word
lengths, $\Sig_k$ is a signed graph, and the seven operons appear as genuine $7$-bundles:
$4919$ of them at $k=31$, organised into chains of which the longest has $698$ vertices. That
chain spans $698+31-1=728$ base pairs, and the corresponding numbers at $k=51$ and $k=75$
give $728$ as well. The agreement with the companion note is literal and not merely
numerical: working on $D_{51}$ rather than on $\Sig_{51}$, and identifying the chain's seven
occurrences against the RefSeq annotation, it reports a longest multiplicity-seven chain of
$s=678$ spanning $728$\,bp whose occurrences hit the seven \texttt{rrn} operons one each.
The quotient neither shortens that chain nor merges it with its own reverse complement, so
$678$ appears again here. On the structural question --- do the five plus-strand and two
minus-strand copies become one seven-copy object in the quotient --- the answer is yes.

On the quantitative question the answer is no, for the two reasons already proved.
Corollary~\ref{cor:nocompute} says $\K(\Sig_k)$ cannot be computed by contraction at this
size, and Theorem~\ref{thm:anti} says that even if it could, it would not count
reconstructions. What is computable, and what is the right residue of the operons, is the
$\mathbb{F}_7$-nullity of the signed Laplacian. The reason it is the right residue is
immediate: modulo $7$ a $7$-bundle contributes $7\equiv0$ to the degree and $\pm7\equiv0$ to
the adjacency, so \emph{every $7$-bundle is invisible over $\mathbb{F}_7$} and every vertex
interior to a $7$-chain yields a zero row. The $\mathbb{F}_7$-nullity therefore measures the
total length of the seven-copy repeat structure, and decays from $7441$ to $5864$ as $k$
rises from $31$ to $75$ and short seven-copy blocks fall below the word length. The
corresponding single-strand figures are $3496$, $3201$ and $2996$: roughly half, as one
expects from a construction that cannot see five copies and two copies as seven.

\section{Scope, provenance, and what is open}

\paragraph{Scope.} Everything proved here concerns the \emph{quotient} $\Sig_k$ and its
undirected shadow. Proposition~\ref{prop:parity}, Lemma~\ref{lem:mono},
Lemma~\ref{lem:switch}, Theorem~\ref{thm:anti} and Theorem~\ref{thm:series} are proved in
full generality. Theorem~\ref{thm:balance} is a combination of Reiner--Tseng
Prop.~11.4--11.5 with the connectivity criterion of the companion note. Theorem~\ref{thm:rt}
is \emph{not} ours: it is Reiner and Tseng's Theorem 1.2, and what is new here is only that
it is exhibited on a reverse-complement de Bruijn graph, including a non-split $2$-primary
extension arising from a real genome. Part (b) of Theorem~\ref{thm:series} --- that the
$r$-rank gain is exactly $s$ rather than merely at least $s$ --- is verified on eight
synthetic pairs $(r,s)$ and on the \emph{E.\ coli} graphs, and is \emph{not} proved.

\paragraph{Provenance.} The signed-graph formalism, the signed incidence and Laplacian
matrices, balance and switching are Zaslavsky's \cite{Zas82}; the dictionary between
bidirected graphs and oriented signed graphs is Edmonds', as recorded in Zaslavsky's glossary
\cite{Zas98}. The critical group of a signed graph and the covering sequence used in
Section~\ref{sec:rt} are Reiner and Tseng's \cite{RT14}; the competing rooted convention and
the chip-firing model for signed graphs are due to Cho, Dochtermann, Inagaki, Oh, Snustad and
Zacovic \cite{CDIOSZ}, building on the invertible-matrix chip-firing of Guzm\'an and Klivans
\cite{GK16}. The spectral form of the double-cover decomposition, which
Theorem~\ref{thm:anti} shows is unavailable here, is Bilu and Linial's \cite{BL06}. The
bidirected-graph model of double-stranded assembly is Medvedev, Georgiou, Myers and Brudno's
\cite{MGMB07}. Sandpile groups of ordinary de Bruijn digraphs are treated by Chan, Hollmann
and Pasechnik \cite{CHP15}. We are not aware of prior work computing critical groups of
reverse-complement de Bruijn graphs, signed or otherwise; a targeted search found none.

\paragraph{Open.} Three things.
\emph{First}, Theorem~\ref{thm:series}(b) in general: prove that a signed $r$-bundle chain of
$s$ internal vertices raises the $r$-rank by exactly $s$, or find a counterexample. A proof
would give a clean partial answer to Question 5.5 of \cite{CDIOSZ}.
\emph{Second}, the series law suggests working with weighted signed graphs over
$\Z[1/(s{+}1)]$, where the contraction does exist; whether the resulting group has any
combinatorial meaning --- whether it counts anything --- is unknown to us.
\emph{Third}, and most important, Theorem~\ref{thm:anti} closes the quotient route to a
double-stranded reconstruction count but does not close the problem. What survives is the
symmetric linking form of Remark~\ref{rem:pairing}, and that is the structure a correct
double-stranded theory should use. Two things are already known about it and both are
constraints rather than encouragements: the form is symmetric, so there are no Lagrangians
in the symplectic sense; and for $\phi$X174 the group order is not a square, so there is no
metabolizer either. A count modulo strand exchange must therefore be a weighted sum over
isotropic subgroups, or over the orbits of the form's isometry group, rather than the order
of a single distinguished subgroup. We have not attempted it, and we note that the naive
guess --- that the answer is $\sqrt{|\K(D_k)|}$ --- is already excluded.

\begin{thebibliography}{99}

\bibitem{BL06} Y.~Bilu and N.~Linial,
\emph{Lifts, discrepancy and nearly optimal spectral gap},
Combinatorica \textbf{26} (2006), no.~5, 495--519.

\bibitem{CHP15} S.~H. Chan, H.~D.~L. Hollmann and D.~V. Pasechnik,
\emph{Sandpile groups of generalized de Bruijn and Kautz graphs and circulant matrices over
finite fields}, J. Algebra (2015), 268--295; arXiv:1405.0113.

\bibitem{CDIOSZ} M.~Cho, A.~Dochtermann, R.~Inagaki, S.~Oh, D.~Snustad and B.~Zacovic,
\emph{Chip-firing and critical groups of signed graphs}, preprint, arXiv:2306.09315v2 (2024).

\bibitem{GK16} J.~Guzm\'an and C.~Klivans,
\emph{Chip firing on general invertible matrices},
SIAM J. Discrete Math. \textbf{30} (2016), no.~2, 1115--1127.

\bibitem{MGMB07} P.~Medvedev, K.~Georgiou, G.~Myers and M.~Brudno,
\emph{Computability of models for sequence assembly},
in: Algorithms in Bioinformatics (WABI 2007), Lecture Notes in Computer Science, Springer,
2007, pp.~289--301.

\bibitem{RT14} V.~Reiner and D.~Tseng,
\emph{Critical groups of covering, voltage and signed graphs},
Discrete Math. \textbf{318} (2014), 10--40; arXiv:1301.2977.

\bibitem{Zas82} T.~Zaslavsky, \emph{Signed graphs},
Discrete Appl. Math. \textbf{4} (1982), no.~1, 47--74; erratum \textbf{5} (1983), 248.

\bibitem{Zas98} T.~Zaslavsky,
\emph{Glossary of signed and gain graphs and allied areas},
Electron. J. Combin., Dynamic Survey \#DS9 (1998).

\end{thebibliography}

\end{document}
